\subsection{集合的基数}

定义1. 若存在从X到Y的双射，则称集合X与集合Y等势。关系“X与Y等势”记为Eq(X, Y)。

关系Eq(X, Y)与Eq(Y, X)显然等价，因此关系Eq(X, Y)是对称的；当它为真时，我们说X与Y是等势的。其次，Eq(X, X)为真。最后，关系Eq(X, Y)是传递的，因为两个双射的复合仍是双射（第二章，§3，第8号，定理1）；因此它是一个等价关系，在每个集合上自反。

由前述可知，若X与Y等势，则关系

\[
(\forall \mathbf {Z}) (\operatorname{Eq} (\mathbf {X}, \mathbf {Z}) \Longleftrightarrow \operatorname{Eq} (\mathbf {Y}, \mathbf {Z}))
\]

是真的。现在，公理模式S7（第一章，§5，第1号）给出了以下公理：

\[
((\forall Z) (\operatorname{Eq} (X, Z) \Longleftrightarrow \operatorname{Eq} (Y, Z)) \Longrightarrow (\tau_ {Z} (\operatorname{Eq} (X, Z)) = \tau_ {Z} (\operatorname{Eq} (Y, Z))).
\]

因此，如果 X 与 Y 等势，我们就有

\[
\tau_ {z} (\operatorname{Eq} (X, Z)) = \tau_ {z} (\operatorname{Eq} (Y, Z)),
\]

这证明了以下定义的合理性：

定义 2. 集合 \(\tau_{\mathbf{Z}}(\operatorname{Eq}(\mathbf{X},\mathbf{Z}))\) 称为 \(\mathbf{X}\) 的基数（或 \(\mathbf{X}\) 的势）并写作 \(\operatorname{Card}(\mathbf{X})\) .

由于 Eq(X, X) 为真，Card(X) 与 X 等势（第一章，§4，方案 S5）。因此我们证明了以下结果：

命题 1. 两个集合 X 和 Y 等势当且仅当它们的基数相等。

\minorheading{示例}

\begin{enumerate}
\def\labelenumi{(\arabic{enumi})}
\item
  \(\operatorname{Card}(\phi)\) 记作 0。由于与 \(\phi\) 等势的唯一集合就是 \(\phi\) （第二章，§3，第 1 和 4 号），我们有 \(0 = \operatorname{Card}(\phi) = \phi\) .
\item
  所有由单个元素组成的集合都是等势的，因为 \(\{(a, b)\}\) 是从 \(\{a\}\) 到 \(\{b\}\) 上的一个双射的图；特别地，它们都与 \(\{\varnothing\}\) 等势。基数 \(\operatorname{Card}(\{\varnothing\}) = \tau_{\mathbf{Z}}(\operatorname{Eq}(\{\varnothing\}, \mathbf{Z}))\) 记为1。 (*)
\item
  \(\operatorname{Card}(\{\varnothing, \{\varnothing\}\})\) 用2表示；这是每一个由两个不同元素组成的集合的基数。
\end{enumerate}

*(4) 可数类型的希尔伯特空间与实数集合等势。*

